
On va
montrer que pour tout @@M0@@, on a :
@@M1@@
Cela entraînera en particulier :
@@M2@@
et montrera donc le théorème.

La démonstration se fait par récurrence descendante sur @@M3@@, le cas @@M4@@ étant trivial. Soit @@M5@@, et posons pour tout @@M6@@ :
@@M7@@
Alors d'après V 6.10,
@@M8@@

On peut donc appliquer l'hypothèse de récurrence, supposant que :
@@M9@@
On a par conséquent :
@@M10@@
D'après 4.4, ceci entraîne :
@@M11@@
Comme @@M12@@ est d'augmentation nulle, @@M13@@.
Comme on fait des calculs modulo torsion, on trouve donc :
@@M14@@

\textbf{Remarque 4.6.1.} On voit que, même si 4.3 était vrai sans tensoriser par @@M15@@, la démonstration ci-dessus ne marche que modulo torsion. Pour un contre-exemple, voir XIV 4.6.

